\clearpage
\VerdierPhysicalPageAnchor{0189}%
\begingroup\small\setlength{\parskip}{0pt}\setlength{\emergencystretch}{1.5em}
\setlength{\abovedisplayskip}{4pt plus 2pt minus 2pt}%
\setlength{\belowdisplayskip}{4pt plus 2pt minus 2pt}%

\VerdierSemanticAnchor{ch3d-2-1-4-proof-cont}
où $v$ est un homotopisme.

Construisons en effet le cône du morphisme $m$ :

\VerdierSemanticAnchor{ch3d-2-1-4-cone-m}
\[
X\xrightarrow{m}Y\xrightarrow{p(m)}c(m)
\longrightarrow X[1],
\]
puis de nouveau, le cône du morphisme $p(m)$ :

\VerdierSemanticAnchor{ch3d-2-1-4-cone-pm}
\[
Y\xrightarrow{p(m)}c(m)\longrightarrow c(p(m))
\xrightarrow{q}Y[1].
\]
% Anomalie diplomatique de l'autorité : « une épimorphisme » ; correction : « un épimorphisme ».
Le morphisme $c(p(m))\xrightarrow{q}Y[1]$ est un épimorphisme
(chap. I, 3.2.2). Il résulte des propriétés des triangles distingués
(propriété TRII) qu'il existe un homotopisme :

\VerdierSemanticAnchor{ch3d-2-1-4-homotopy-equivalence}
\[
X\xrightarrow{v}c(p(m))[-1]
\]
tel que le triangle ci-après :

\VerdierSemanticAnchor{ch3d-2-1-4-triangle}
\begin{center}
\begin{tikzcd}[column sep=3.5em,row sep=3.2em]
X \arrow[rr,"v"] \arrow[dr,"m"']
  && c(p(m))[-1] \arrow[dl,"{q[-1]}"] \\
& Y &
\end{tikzcd}
\end{center}
soit commutatif à homotopie près.

On peut donc supposer que le morphisme $m$ est un épimorphisme de
$\operatorname{comp}(\mathcal A)$. Soit alors :

\VerdierSemanticAnchor{ch3d-2-1-4-kernel}
\[
N=\cdots\longrightarrow N^i\xrightarrow{d_N^i}N^{i+1}
\xrightarrow{d_N^{i+1}}N^{i+2}\longrightarrow\cdots
\]
le noyau du morphisme $m$. Posons :

\VerdierSemanticAnchor{ch3d-2-1-4-truncation}
\[
N_{-n}=\cdots\longrightarrow N^{n-3}\longrightarrow N^{n-2}
\longrightarrow\operatorname{Im}(d_N^{n-2})
\longrightarrow0\longrightarrow0\longrightarrow\cdots.
\]
Les $N_{-n}$ forment une suite décroissante de sous-complexes acycliques
de $N$, d'où, en posant $X_n=X/N_n$, une suite :

\VerdierSemanticAnchor{ch3d-2-1-4-quotients}
\[
u_n\colon X_n\longrightarrow X_{n-1},\qquad n\in\mathbb Z.
\]
Les $u_n$ sont des morphismes élémentaires de hauteur $-n$ et
\[
\varprojlim_n X_n=X,
\qquad
\varinjlim_n X_n=Y.
\]
\endgroup
